\scanpage{58}

Considérons le composé
\[
\Delta^*j:\quad\underline{\operatorname{NS}}_{A/S}\deleted{(\unsure{A})}\longrightarrow P\deleted{\underline{\unsure{\mathrm{ic}}}}_{A/S}
\]
on \unsure{sait|voit} que son composé avec $P\unsure{\deleted{\underline{\mathrm{ic}}}}_{A/S}\overset{k}{\unsure{\longrightarrow|\mathbin{-}}}\underline{\operatorname{NS}}_{A/S}$ est $2\operatorname{id}_{\underline{\operatorname{NS}}_{A/S}}$. Donc \uline{le $\overset{\inserted{\unsure{2|e}}}{\deleted{\unsure{e|2}}}$ double de l'extension canonique de} $\underline{\operatorname{NS}}_{A/S}$ \uline{par} $P_{A/S}$ \uline{splitte}

\begingroup
\setlength{\unitlength}{\dimexpr\linewidth/12\relax}
\noindent\begin{picture}(12,2.10)
\put(0,0){\begin{tikzpicture}[x=\unitlength,y=\unitlength,baseline=0pt]
\path[use as bounding box] (0,0) rectangle (12,2.10);
\draw[line width=.4pt] (0.20,0.30)
  .. controls (0.80,0.20) and (2.30,0.34) .. (2.70,0.59)
  .. controls (2.95,0.73) and (2.76,1.14) .. (2.79,1.41)
  .. controls (2.81,1.60) and (3.17,1.63) .. (3.45,1.61)
  .. controls (5.50,1.60) and (9.70,1.69) .. (11.00,1.75);
\end{tikzpicture}}
\put(0,1.10){\uline{canoniquement}.}
\put(3.68,1.10){Considérons d'autre part}
\end{picture}\par
\endgroup

\noindent la restriction de $j'(k\Delta^*)\deleted{\unsure{g|j}}$ à $\underline{\operatorname{Tors}\,\unsure{\mathrm{birig}}\,\mathrm{sym}}(A,A;G)$.
Je dis que c'est la multiplication par $2$\unsure{.}, i.e.
\[
\mathord{\hbox{\rlap{\begin{tikzpicture}[overlay,baseline=0pt,x=1ex,y=1ex]
\draw[line width=.3pt] (-.65,.72) -- (-.61,1.05);
\draw[line width=.3pt] (-.70,-.08) .. controls (-.59,-.12) and (-.54,-.02) .. (-.52,.10);
\end{tikzpicture}}}}k\Delta^*
\]

\margin{57}
